\documentclass[12pt]{article}
\usepackage{amsmath}
\usepackage{amsthm}
\usepackage{bm}
\usepackage[hmargin=1in,vmargin=1in]{geometry}
\usepackage[shortlabels]{enumitem}
\usepackage{multicol}
\usepackage{hyperref}
\usepackage{graphicx}
\usepackage[section]{placeins}
\usepackage{amssymb}
\usepackage{cite}
\usepackage{listings}
\usepackage{amsmath,tkz-euclide}
\usepackage{tikz}
\usetikzlibrary{arrows.meta, decorations.markings}
\usetkzobj{all}
\counterwithin{figure}{section}
\graphicspath{{/Users/edwardmcdugald/NaturalStripePatterns/presentations/figs/}}
\numberwithin{equation}{section}

\parindent 0in
\parskip 1em
\Urlmuskip=0mu plus 1mu

\title{Computational Methods and Results Concerning the RCN Approximation to Swift-Hohenberg Fields}
\author{Edward McDugald}

\begin{document}

\maketitle

\section{Introduction}
%Describe ubiquitousness and examples of pattern forming systems
%Describe SH as a canonical pattern forming PDE
%Mention that there is a theory based on perturbation expansion that captures a range of defects
%Mention that the report will cover analytical derivation of RCN, computational verification of RCN, and systems identification attemps
\par Defect ridden striped patterns are ubiquitous in nature, evident in zebra stripes, sand ripples, convecting fluids, and wideband lasers. The canonical microscopic model for such systems is the 2D Swift-Hohenberg (SH) PDE. In this model, one may interpret the microscopic variable as a height averaged temperature or vertical component of velocity for a convecting fluid. The character of the long term field configuration is highly dependent on the domain, boundary conditions, and initial condition. That is to say, these variables can give rise to a range of topologically distinct defect types. Regardless of what types of defects are present, the field will locally resemble a plane wave almost everywhere. Indeed, on infinite domains, the SH system has stationary solutions of the form $\sum_{n}A_n(k^2)\cos{n\theta}$, where $\vec{k} = \nabla \theta$.
\par The theory of natural stripe patterns is largely focused on obtaining macroscopic order parameter equations for SH systems that are solved on finite domains. In the finite domain setting, one may introduce a small parameter $\epsilon = \frac{l}{L}$, where $l$ is the height of the container, and $L$ is maximum width of the container. Using this small parameter, and the assumption that locally the field resembles a plane wave, one can introduce multiple length scales in conjunction with a phase and amplitude ansatz to yield phase and amplitude equations via a perturbation expansion. Equivalently, one may average over the periodicity of the field to yield the same macroscopic equations.
\par The leading order phase equation is known as the Cross-Newell (CN) phase diffusion equation. While solutions of CN capture topological properties of SH defects, they allow for a much wider range of spatial wave numbers compared to SH solutions. To second order, one obtains the regularized Cross-Newell (RCN) equation, which captures both the topological properties of defects and the narrow range of wave numbers seen in SH simulations. The amplitude equation is algebraic in $k$ to leading order. To second order, it is a proper PDE, however the second order correction is only important near defects. Thus, the hope is that the second order phase equation, in conjunction with the leading order algebraic amplitude equation, yeild good approximations to SH fields away from defects.
\par The RCN may be written as a gradient flow, from which stationary field configurations can be sought by seeking minimizers of the energy. One strategy is to seek self-dual minimizers. Such minimizers yield PDEs that can be solved analytically, and provide closed form formulas for phase grain boundaries and stationary dislocations. They can also be related to solutions of the Helmholtz equation to deduce canonical defects such as concave and convex disclinations, saddles and targets. While much analysis has been done on the RCN equation and its self dual solutions, there are many open challenges. For instance, a particular RCN solution may give a good formula for a particular defect in isolation, however analysis is lacking in situations where multiple defects are interacting. For instance, SH solutions on the ellipse appear to reflect a competition between the eikonal solution, phase grain boundaries, and dislocations.
\par In this report, we will go over the formal derivation of the RCN equations and discuss examples of self-dual RCN solutions that yield qualitative agreement to SH fields. We will then discuss efforts to fit these RCN solutions to actual SH simulations. Finally, we will discuss systems identification, and propose SH as an interesting test case for SINDy-like methods that can deduce order parameter equations from microscopic data.


\begin{figure}[h!]
\centering
\includegraphics[width=0.45\textwidth]{ferrofluid.png}
\includegraphics[width=0.45\textwidth]{optical_interference.png} \\[2ex]
\includegraphics[width=0.45\textwidth]{sandripples.png}
\includegraphics[width=0.45\textwidth]{zebra.png}
%\includegraphics[width=0.45\textwidth]{vegetation.png}
\label{fig:natural patterns}
\caption{Striped patterns in nature. Ferrofluid, optical interference, sand ripples, vegetation, zebra stripes}
\end{figure}

\section{RCN Derivation and analysis}
\par The procedure for deriving such a phase equation involves studying the vector field whose components are the instantaneous frequencies of the pattern generated by the microscopic equations. We call this vector field $\bm{k}$, and note that this is the gradient of the phase, $\bm{k} = \nabla \theta$. The phase gradient is modulated over distances $L$, the average distance between defects, which is long with respect to the local wavelength $l$. We introduce the following variables:
\begin{align*}
\bm{X} &= \epsilon\bm{x}\\
T &= \epsilon^2 t\\
\Theta(\bm{X},T) &= \epsilon \theta(\bm{x},t)\\
k &= \|\bm{k}\|\\
\epsilon &= \frac{l}{L} \quad \text{(The inverse aspect ratio)}.
\end{align*}
Letting $w$ be the variable of the microscopic equations, and $E$ the associated energy functional, we apply the following modulational Ansatz:
\begin{align}
w &= \sum A_n(k^2)\cos(n\theta) \quad \text{(for real fields)}\\
w &= A(k^2,\bm{X},T)e^{i\theta(\bm{x},t)} \quad \text{(for complex fields)}.
\end{align}
We then average the energy over the periodicity of the pattern
\begin{equation}
\overline{E} = \frac{1}{2\pi} \int E d\theta
\end{equation}
and analyze the dependence of the energy on $A$ and $\Theta$ in orders of $\epsilon$. We find that the only orders of $\epsilon$ that contribute are $\mathcal{O}(\epsilon^0)$ and $\mathcal{O}(\epsilon^2)$, and that the energy can be interpreted as a sum of stretching and bending energies. Figure 2.1 provides a heuristic interpretation for the macroscopic model, where macroscopic energy densities are estimated from microscopic energy densities by convolution with a Gaussian.
\begin{figure}[h!]
\centering
%\includegraphics[width=0.8\textwidth]{emacro_est_fft2.png}
\includegraphics[width=0.8\textwidth]{emacro_est_fft3.png}
\label{fig:micro vs macro energy}
\caption{Depiction of Microscopic vs. Macroscopic energy. The top figures correspond to microscopic energy densities, the bottom figures correspond to macropscopic energy densities.}
\end{figure}
\FloatBarrier

\subsection{Derivation from Complexified Swift Hohenberg}
\par In its complexified form, the swift hohenberg equation has energy density
\begin{equation}
E = \int \left( (\nabla^2+1)w(\nabla^2+1)w^{*}-Rww^{*}+\frac{1}{2}w^2w^{*^2}\right)dx.
\end{equation}
Let us compute $\delta E$,
\begin{align*}
\delta E &= \lim_{h \rightarrow 0}\frac{E[w+h\delta w, w^* + h\delta w^*]-E[w,w^*]}{h}\\
    &= \begin{aligned}[t]
        & \int \left( (\nabla^2+1)w(\nabla^2+1)\delta w^* -Rw\delta w^* + w^2w^*\delta w^* + \right. \\
        & \left. (\nabla^2+1)\delta w (\nabla^2+1)w^* -Rw^*\delta w + ww^{*^2}\delta w \right) d\bm{x}
    \end{aligned}
\end{align*}
Using integration by parts, one finds that
\begin{align*}
(\nabla^2+1)w(\nabla^2+1)\delta w^* &= (\nabla^2 +1)^2w\delta w^*\\
(\nabla^2+1)\delta w (\nabla^2+1)^2w^* &= (\nabla^2+1)^2w^*\delta w.
\end{align*}
Thus, one obtains
\begin{align}
\frac{\delta E}{\delta w^*} &= (\nabla^2+1)^2w - Rw + w^2w^*\\
\frac{\delta E}{\delta w} &= (\nabla^2+1)^2 w^* -Rw^* + w^{*^2}w.
\end{align}
The complexified Swift-Hohenberg equations can be written as
\begin{equation}
w_t = - \frac{\delta E}{\delta w^*}, \quad w^*_t = -\frac{\delta E}{\delta w}.
\end{equation}
Using the Ansatz for complex fields (1.2), one can write
\begin{equation}
\delta E = \frac{\delta E}{\delta w}\delta w + \frac{\delta E}{\delta w^*}\delta w^* = -2A_t\delta A - 2A^2\Theta_t \delta \theta.
\end{equation}
\par We want now to find expressions for $\frac{\delta E}{\delta A}$ and $\frac{\delta E}{\delta \Theta}$.
Using the fact that $\nabla_{\bm{x}} = \epsilon\nabla_{\bm{X}}$, and writing $w=Ae^{i\theta}=A(\bm{x})e^{i\theta(\bm{x})}$, one obtains
\begin{align}
\nabla_{\bm{x}}w &= e^{i\theta}(i\bm{k}+\epsilon\nabla_{\bm{X}})A\\
\nabla^2_{\bm{x}}w &= e^{i\theta}(-k^2+i\epsilon(2\bm{x}\cdot \nabla_{\bm{X}}+\nabla_{\bm{X}}\cdot\bm{k})+\epsilon^2\nabla^2_{\bm{X}})A\\
\nabla_{\bm{x}}w^* &= e^{-i\theta}(-i\bm{k}+\epsilon\nabla_{\bm{X}})A\\
\nabla^2_{\bm{x}}w^* &= e^{-i\theta}(-k^2-i\epsilon(2\bm{x}\cdot \nabla_{\bm{X}}+\nabla_{\bm{X}}\cdot\bm{k})+\epsilon^2\nabla^2_{\bm{X}})A.
\end{align}
Inserting (1.2) into (2.1) with (2.6)-(2.9), and \emph{ignoring} terms of order $> \mathcal{O}(\epsilon^2)$, one obtains
\begin{align}
E &= \begin{aligned}[t]
     &\int \left( (k^2-1)^2A^2 - RA^2 + \frac{1}{4}A^4 + K\right)d\bm{x}\\
     & + \epsilon^2 \int \left( (2\bm{k}\cdot \nabla A + \nabla \cdot \bm{k}A)^2 + 2(1-k^2)A\nabla^2A \right)d\bm{x}.
     \end{aligned}
\end{align}
Now, obtaining macroscopic equations in terms of the phase parameter requires averaging the energy over the periodicity of the pattern. Conveniently, in the complex field case, this averaging occurs automatically, and we have $E = \overline{E}$. Equation (2.5) gives
\begin{equation}
-2A_t = \frac{\delta \overline{E}}{\delta A} = \frac{\delta\overline{E_0}}{\delta A}+\epsilon^2 \frac{\delta \overline{E_2}}{\delta A}.
\end{equation}
By construction, $A_t = \epsilon^2 A_T$, and therefore $A_t$ is at most $\mathcal{O}(\epsilon^2)$. Thus, to leading order, the main equation contributing to the evolution of the amplitude $A$ is
\begin{equation}
\frac{\delta \overline{E_0}}{\delta A} = 0.
\end{equation}
This gives
\begin{equation}
A^2 = R - (k^2-1)^2.
\end{equation}
This is known as the eikonal solution, which demands that the amplitude be slaved algebraically to the modulus of the phase gradient. This property holds for stress parameter values far from onset.
\par Having found a dominant relation for $A$, we seek an equation for $\Theta$. Using the eikonal solution, we can write the leading order energy as
\begin{equation}
\overline{E_0} = \int - \frac{1}{2}\left( A^2(k^2) \right)^2 + K.
\end{equation}
Considering just the $\mathcal{O}(\epsilon^0)$ contribution, equation (2.5) gives
\begin{equation}
A^2 \theta_t = -\frac{1}{2}\frac{\delta \overline{E_0}}{\delta \theta}.
\end{equation}
We have
\begin{align*}
\frac{\delta \overline{E_0}}{\delta \theta}\delta \theta &= \delta\left( -\frac{1}{2}(A^2(k^2))^2\right)\\
                                                         &= -A^2(k^2)\frac{d}{dk^2}A^2(k^2)2\bm{k}\cdot \delta \nabla \theta.
\end{align*}
Using integration by parts, one finds that
\begin{equation*}
-A^2(k^2)\frac{d}{dk^2}A^2(k^2)2\bm{k}\cdot \delta \nabla \theta = A^2(k^2)\frac{d}{dk^2}A^2(k^2)2\nabla \cdot \bm{k}\delta \theta.
\end{equation*}
Thus, we obtain
\begin{equation*}
A^2 \theta_t = -A^2(k^2)\frac{dA^2(k^2)}{dk^2}2\nabla\cdot\bm{k}
\end{equation*}
Noting that $\theta_t = \epsilon \Theta_T$, $\nabla_{\bm{x}}=\epsilon\nabla_{\bm{X}}$, and setting
\begin{equation}
B(k^2) = A^2(k^2)\frac{dA^2(k^2)}{dk^2},
\end{equation}
one obtains the \emph{unregularized} Cross-Newell phase diffusion equation
\begin{equation}
A^2\Theta_T + \nabla_{\bm{X}}\cdot \bm{k}B(k^2) = 0.
\end{equation}

Recall the form of $\overline{E_0}$, $\overline{E_0} = \int -\frac{1}{2}A^4(k^2)+K$.
Let us define $k_B$ to be the preferred wave number, by which we mean the wave number that minimizes $E_0$.
That is, $k_B$ is defined to satisfy the equation
\begin{equation}
\frac{d}{dk^2}A^4(k^2) = 0.
\end{equation}
We define $K$ to be $-\frac{1}{2}A^4(k^2)$ evaluated at $k_B$, and we can then write
\begin{equation}
\overline{E_0} = \int \left( -\frac{1}{2}A^4(k^2)\right)_{k^2}^{k_B^2}d\bm{x}
\end{equation}
Note that this choice of $K$ ensures that the leading order energy is $0$ whenever $k = k_B$.
\par Let us look at the unregularized equation (2.17). We can write this as a diffusion equation in standard form
\begin{equation}
A^2\Theta_T = -\nabla_{\bm{X}}\left[ B(k^2)\nabla \theta\right].
\end{equation}
Expanding the right hand side yields
\begin{equation}
A^2\Theta_T = \theta_{xx}\left(-\frac{dB(k^2)}{dk^2}2\theta_x^2 - B(k^2)\right) + \theta_{xy}\left(-\frac{dB(k^2)}{dk^2}4\theta_x\theta_y\right) + \theta_{yy}\left(-\frac{dB(k^2)}{dk^2}2\theta_y^2 -B(k^2)\right).
\end{equation}
Consider the simple case where $\bm{k} = \langle \theta_x, 0\rangle$.
We then obtain
\begin{equation*}
A^2\Theta_T = \theta_{xx}\left(-\frac{dB(k^2)}{dk^2}2\theta_x^2 - B(k^2)\right) + \theta_{yy}\left(-B(k^2)\right).
\end{equation*}
To have a well-posed diffusion equation, we need that
\begin{align}
-\frac{dB(k^2)}{dk^2}2\theta_x^2 - B(k^2) &>0\\
-B(k^2) &>0
\end{align}
Noting that
\begin{equation}
\frac{d}{dk}\left(kB(k^2)\right) = B(k^2)+k\frac{d}{dk}B(k^2) = B(k^2)+2k^2\frac{d}{dk^2}B(k^2),
\end{equation}
We see that we require
\begin{align}
\frac{d}{dk}(kB(k^2)) &<0\\
B(k^2) &< 0.
\end{align}
It turns out that $B, (kB)'$ are the eigenvalues of the second order quasi-linear operator, and for all $k$ the above must hold. Let us consider the plot of $k$ vs $kB$. For the case of Swift-Hohenberg with $R=.5$, we have $B = -4(k^2-1)(.5-(k^2-1)^2)$. We see that there is a limited range of wave numbers where these values are satisfied. This range of wave numbers is known as the Busse-Balloon, and is the range $k_B<k<k_E$ as labeled in figure 2.2. It turns out for natural patterns, almost everywhere $k=k_B$, however along line defects we will have $k<k_B$.
%\begin{figure}
%\centering
%\includegraphics[scale=0.5]{mybb.png}
%\caption{Busse Balloon for Swift Hohenberg, $R=.5$}
%\end{figure}
\begin{figure}
\centering
\includegraphics[scale=0.5]{KB.png}
\caption{Generic Busse Balloon}
\end{figure}

Consider, for example, the scenario depicted in figure 2.3. We see that along the line defect, the pattern has a longer wavelength. This means that the local spatial frequences are smaller than they are along most of the pattern, where the wavenumber is close to $k_B$. Thus, along the line defect $k<K_B$, which lies outside of the Busse Balloon. Hence, the unregularized Cross-Newell phase diffusion equation is ill-posed along line defects, and regularization is required.
\begin{figure}
\centering
\includegraphics[scale=2.0]{wavenumber_mismatch.png}
\caption{Change in wavelength along line defects}
\end{figure}
\par To regularize the Cross-Newell phase diffusion equation, we analyze the $\mathcal{O}(\epsilon^2)$ part of the energy, $\overline{E_2}$.
We have that
\begin{equation}
\overline{E_2} = \int \left((2\bm{k}\cdot \nabla A+\nabla \cdot \bm{k}A)^2+2(1-k^2)A\nabla^2A\right)d\bm{x}.
\end{equation}
It turns out that since $\nabla_{\bm{x}} = \bm{k}\frac{\partial}{\partial \theta}+\epsilon \nabla_{\bm{X}}$, the terms either involve differentiating $\bm{k}$ as it changes direction along the phase contour, or differentiation of $k^2$ in the direction of $\bm{k}$. The terms $2\bm{k}\cdot \nabla A$ and $A\nabla^2A$ are small, and $\overline{E_2}$ can be well approximated by
\begin{equation}
\int A^2(k_B^2)(\nabla \cdot \bm{k})^2d\bm{x}.
\end{equation}
Computing its variation, we obtain the regularized Cross-Newell phase diffusion equation
\begin{equation}
A^2(k_B^2)\Theta_T + \nabla \cdot \bm{k}B + \epsilon^2 A^2(k_B^2)\nabla^4\Theta = 0.
\end{equation}
Moreover, we can approximate the energy $\overline{E}$ by
\begin{equation}
\overline{E} = \int \left(\left( (\frac{d}{dk^2}B(k^2))\right)_{k_B^2}(k^2-k^2_B) + \epsilon^2 (\nabla \cdot \bm{k})^2A^2(k_B^2)\right)d\bm{x}.
\end{equation}
The $\mathcal{O}(\epsilon^0)$ contribution is the "strain" energy, and the $\mathcal{O}(\epsilon^2)$ contribution is the "bending" energy.

\subsection{RCN for Real Fields}
In the case of real fields, the Swift-Hohenberg equation has microscopic energy functional
\begin{equation}
E = \int \left( ((\nabla^2+1)w)^2-\frac{1}{2}Rw^2 + \frac{1}{4}w^4 \right)d\bm{x}.
\end{equation}
In this case, the microscopic equation is
\begin{equation}
w_t = -\frac{\delta E}{\delta w} = -(\nabla^2+1)^2w + Rw - w^3.
\end{equation}
Note that equation (3.2) admits stationary solutions for striped patterns
\begin{equation}
w(\bm{x},t) = w(\theta; \{A_n(k^2)\},R) = \sum A_n(k^2)\cos(n\theta).
\end{equation}
In the real field case, our modulation ansatz is given by (1.1),
\begin{equation}
w(\bm{x},t) = w(\theta; \{A_n(k^2)\},R; \nabla \theta = \bm{k}(\bm{X},T)=\nabla_{\bm{X}}\Theta(\bm{X},t)).
\end{equation}
The procedure to derive the phase diffusion equation is the same as it was for the case of complex fields. However, the energy is not automatically averaged, making the procedure far more delicate. We simply state the key results, which are derived by Newell \cite{NEWELL2023}.
\par
The microscopic equations evolve according to
\begin{equation}
 w_t \delta w = -\delta E.
\end{equation}
Upon averaging over $\theta$, one obtains
\begin{equation}
 \langle w_{\theta}^2\rangle\theta_t \delta \theta = -\delta \overline{E}
\end{equation}
The averaged energy to leading order is given by
\begin{equation}
\overline{E_0} = \left(-\frac{1}{4}\langle w^4 \rangle + K \right)dxdy
\end{equation}
Again, $k^2_B$ is defined so as to satisfy
\begin{equation}
\frac{d}{dk^2}\left(\frac{1}{4}\langle w^4 \rangle\right) = 0,
\end{equation}
and defining $K$ as we did for the complex case, obtain
\begin{equation}
\overline{E_0} = \int \left( \frac{1}{4}\langle w^4 \rangle\right)_{k^2}^{k^2_B}dxdy.
\end{equation}
Considering just the leading order contribution of the energy, we can obtain the unregularized diffusion equation by taking the variaton of $\overline{E_0}$, and obtain
\begin{equation}
\langle w_{\theta}^2 \rangle \Theta_T + \nabla_{\bm{X}}\bm{k}B(k^2)=0,
\end{equation}
where
\begin{equation}
B(k^2) = 2\frac{dE_0}{dk^2},
\end{equation}
where $E_0$ here is used to denote the integrand of $\overline{E_0}$. Note that in the complex case, we had $B(k^2)=\frac{dE_0}{dk^2}$.
\par The fact that the $\mathcal{O}(\epsilon)$ terms vanish in the averaged energy is not as clear in the case of the real field, but it still holds. The argument as to why is omitted for the time being.
\par Considering the $\mathcal{O}(\epsilon^2)$ terms, we argue in a similar fashion as we did for the case of complex fields. All terms involving $(\bm{k}\cdot \nabla_{\bm{X}})$ acting on the amplitudes yield terms proportional to $\frac{dA_n(k^2)}{dk^2}(\bm{k}\cdot \nabla_{\bm{X}}k^2)$ which can be written as $\bm{k}\cdot \nabla_{\bm{X}}(k^2-k_B^2)$. Since $k^2-k_B^2$ is small almost everywhere, these terms are negligible compared to the terms proportional to $\nabla \cdot \bm{k}$. Via integration by parts, one can obtain
\begin{equation}
\epsilon^2 \int \langle w_{\theta}^2 \rangle (\nabla \cdot \bm{k})^2 d\bm{x} \approx \epsilon^2 \langle w_{\theta}^2 \rangle_{k_B}\int (\nabla \cdot \bm{k})^2 d\bm{x}.
\end{equation}
The regularized phase diffusion equation is then
\begin{equation}
\langle w_{\theta}^2 \rangle \Theta_T + \nabla \cdot \bm{k}B + \epsilon^2 \langle w_{\theta}^2 \rangle\nabla^4 \Theta = 0.
\end{equation}

\subsection{Analysis of RCN and its Solutions}
As we did for the complex field case, we can write the phase diffusion equation as a gradient flow, with the form
\begin{align}
\eta \nabla^4 \theta + \nabla \cdot \bm{k}B &= -\frac{\delta \overline{E}}{\delta \theta}\\
\overline{E} &= \int \left( \frac{1}{2}\eta|\nabla^2\theta|^2 + \frac{\alpha}{2}G^2 \right)d\bm{x}\\
G^2 &= -\frac{1}{\alpha}\int_{k_B^2}^{k^2}Bdk^2\\
\alpha &= -\frac{1}{2}\frac{dB}{dk^2}\vert_{k^2=1}.
\end{align}
In this form, the first term measures the "bending" energy, whereas the second term measures the "stretching" energy. Figure 2.4 shows a solution of Swift-Hohenberg on the square alongside its energy density. We see that the energy tends to condense around bends within the pattern, and that the final state has an almost constant spatial wavenumber. We use random initial conditions for this solution.
\begin{figure}
\centering
\includegraphics[scale=0.95]{sh_eq_2.png}
\caption{Swift-Hohenberg solution on the square, $R=.5$}
\end{figure}

\par Motivated by observations of the stationary dislocation \cite{NEWELL1996}, a key component of the RCN analysis to date is to seek minimizers of the energy that are self-dual or anti self-dual. That is, we seek solutions of the equation
\begin{equation}
    \frac{1}{2}\eta|\nabla^2 \theta|^2 = \frac{\alpha}{2}G^2.
\end{equation}
To this end, let us analyze
\begin{equation}
    \nabla^2\theta = \beta s G + s \chi,
\end{equation}
where
\begin{equation}
    \beta = \sqrt{\frac{\alpha}{\eta}}, s = \pm 1.
\end{equation}
After some manipulation, one finds
\begin{equation}
    \nabla^2\chi + \nabla s\chi \nabla_{\vec{k}}G = -\beta s \nabla_{\vec{k}}\cdot \left(\hat{J}(\hat{I}-\nabla^2\theta I)\nabla_{\vec{k}}G\right).
\end{equation}
The $s$ terms captures the sign ambiguity in the self-dual balance, and $\chi$ measures the deviation from self-duality.
The term $G^2$ captures the stretching component of the macroscopic energy. By construction, $\alpha$ is defined so that to loeading order, $G^2 = (k^2 - 1)^2.$ (2.52) may then be well approximated by
\begin{equation}
    \nabla^2\theta = \beta s (1-k^2)+s\chi.
\end{equation}
Self-dual solutions may then be well approximated by solutions of
\begin{equation}
    \nabla^2\theta = \beta s (1-k^2)+s\chi.
\end{equation}
Letting $\theta = \frac{1}{\beta s}\ln{\psi}$ and $\chi = \psi v$, one obtains the PDE pair
\begin{align}
    \nabla^2 \psi - (\beta^2 + \beta\chi)\psi &= 0\\
    \nabla^2v - (\beta^2+\beta\chi)v &= -4\beta \psi^{-1}J.
\end{align}
An important aspect of (2.58) is that the right hand side has a factor of $\hat{J}\vec{k}$. It turns out that if $k^2 = 1$, $\vec{k}$ is an eigenvector of $\hat{J}$ with eigenvalue $0$. Thus, over most of the pattern domain/away from defects, $\chi = 0$.
Thus, when $J = 0$, $\chi = 0$, and we are simply left with the helmholtz equation
\begin{equation}
    \nabla^2 \psi - \beta^2\psi = 0.
\end{equation}
This gives rise to the simplest self-dual solutions, which are sums of exponentials,
\begin{equation}
    \psi = \sum_{1}^{N}\exp{\beta \vec{k}_j \cdot x}, |\vec{k}_j| = 1.
\end{equation}
The far fields of phase grain boundaries, concave and convex disclinations, targets, saddles, and spines can be realized through such solutions.

\begin{figure}
\centering
\includegraphics[scale=0.7]{concaveconvex.png}
\caption{Concave/Convex disclinations}
\end{figure}
\FloatBarrier

\begin{figure}
\centering
\includegraphics[scale=0.7]{spine.png}
\caption{Spine Defect}
\end{figure}

\FloatBarrier


\par Another interesting defect realized from RCN is the stationary dislocation. Such defects are obtained by seeking solutions $\psi$ such that $\psi = \exp{\beta x}F(x,y)$, where $2\beta F_{x} + F_{yy} = 0$. We seek solutions such that $\theta(-x,y)=\theta(x,y)$ and that $\theta(x,y) = x$ for $y$ large and positive and to $x+\pi \text{sgn}(x)$ for $y$ large and negative. The solution is then
\begin{equation}
    \theta(x,y) = x + \frac{1}{\beta}\ln\{\frac{1}{2}\left(1+\exp{(\beta \pi \text{sgn}(x))}\right) + \frac{1}{2}(1-\exp{(\beta \pi \text{sgn}(x))})\text{Erf}(\sqrt{\beta}\frac{y}{\sqrt{|x|}})\}.
\end{equation}


\subsection{Open Challenges for the RCN Model}
\subsubsection{SH Zippers and the Handle Formation}
\par The case of $N=2$ describes a prevalent defect known as a phase grain boundary.
The phase field is given by the formula
\begin{equation}
\theta_{\mu}(x,y) = \sqrt{1-\mu^2}x + \log{2\cosh{\mu y}}.
\end{equation}
While this solution arises from the self-dual reduction, this solution is indeed the minimizer of the RCN energy within a function space of single-valued phases. This is valid for arbitrary $\mu \in (0,1)$. However, numerical solutions to Swift-Hohenberg do not exhibit phase grain boundaries past a certain critical angle. Instead, dislocations are formed along the core of the defect, as depicted in figure 2.7.
\par The difference between phase grain boundaries and dislocations are well understand by comparing properties of the phase and its gradient. For small $\mu$, the $y$ component of the wave vector $\theta_y$ vanishes at the defect core. As $\mu$ increases and a dislocation forms, a jump is developed across the defect core. In addition, along the $x$ direction the phase alternates between regions where it is constant, and regions where it is linear with $\theta_y$ constant. These suggest mixed boundary conditions for the phase at the defect core.
\par The most natural question to ask is whether the RCN equation predicts this transition of boundary conditions for the phase as $\mu$ is increased. It is known that regardless of $\mu$, phase grain boundaries are minimizers of the RCN energy when restricted to a function space of single valued phases. However, it is also known that if this requirement is relaxed, there exist lower energy states in the space of two valued phases. In fact, numerical analogues of the SH zippers have been obtained for the RCN equations \cite{VENK2008}, as seen in figure 2.8.
\begin{figure}
\centering
\includegraphics[scale=0.4]{sh_zippers.png}
\caption{Swift-Hohenberg Zippers}
\end{figure}
\begin{figure}
\centering
\includegraphics[scale=0.5]{CN_Zip.png}
\caption{Cross-Newell Zippers}
\end{figure}






\subsubsection{Convection in an Elliptical Container}
\par While there is ample evidence that the RCN model captures the transition observed from a phase grain boundary to a zipper, the scenario of convection rolls in an elliptical container remains an open challenge. That is, there is no formal argument to explain the range of character of defects observed in this scenario, nor is there a numerical solution of RCN yielding a qualitatively similar field.
\par Figure 2.9 shows a nearly stationary SH solution and its corresponding energy density on the ellopse. Using the eikonal solution as an initial condition, one observes the formation of dislocations along the major axis of the ellipse, reminiscent of the zippers observed in the sharp knee bend scenario. However, as one moves away from the center of the ellipse to its boundary in the $x$ direction, one sees the dislocations soften into phase grain boundaries. In addition, the phase contours appear eikonal-like at the outer most layer, and gradually develop a more jagged profile moving into the center of the ellipse. Due to the finite domain and more restrictive boundary conditions, performing a variational analysis on this scenario is more delicate than the case of the zippers. Similarly, numerical solutions are more delicate to obtain.
\begin{figure}
\centering
\includegraphics[scale=0.8]{SHEllipse2to1_0124.png}
\caption{SH Solution on the Ellipse}
\end{figure}


%Go to into derivation for complex fields
%mention that it can be done for real fields
%Show formula for various defects given from self dual reduction

\section{Computational Verification}
\par By construction, the RCN model is valid under the assumption of a slowly varying wave vector. Thus, near defects, one doesn't expect agreement between RCN solutions and Swift-Hohenberg solutions. Remarkably, however, the RCN solutions valid away from defects preserve the topological characteristics of defects. In this section, we explore the accuracy of the RCN model in the case of the phase grain boundary and the stationary dislocation. We compare the result of numerical swift-hohenberg fields to patterns obtained using RCN phase formula, and the amplitude relations
\begin{align}
    A_1 &= \frac{4}{3}(R - (k^2-1)^2)\\
    A_3 &= \frac{A_1^3}{4(R-(9k^2-1)^2)}.
\end{align}

\par The general procedure is as follows. We simulate swift-hohenberg numerically with desired defect type. From this simulation, we use an optimization procedure to fit parameters of the RCN phase function to the Swift-Hohenberg data, and compute the error. With potential systems identification applications in mind, we make it possible to estimate initial values for the RCN parameters from the SH data. The data is simulated using a fourth order solver for stiff PDEs introduced by Trefethen \cite{TREF2005}. Depending on the desired defect, boundary conditions can be modified by ramping the stress parameter $R$ over an appropriate domain. 
\subsection{Phase Grain Boundary}
In the case of a PGB, one can write the phase as
\begin{equation}
    \theta = \log(\exp(k_1 x + k_2 y)+\exp(k_1 x - k_2 y)).
\end{equation}
Thus, the parameters we are optimizing are $k_1, k_2$. We also allow for a phase shift parameter, $\phi$.
The amplitude equations are given by (3.1) and (3.2).
Letting $w$ denote our numerically obtained swift-hohenberg field, we seek parameters $k_1,k_2,\phi$ by minimizing the objective function
\begin{align}
    \text{obj}(k_1,k_2,\phi) = \left( A_1\cos(\theta(k_1,k_2)-\phi) + A_3\cos(3(\theta(k_1,k_2)-\phi))- w\right)^2.
\end{align}
We terminate the optimization procedure once the gradient of the objective is smaller than $0.00001$. The recovered parameters are
\begin{align*}
    k_1 &= .952\\
    k_2 &= .300\\
    \phi &= -1.532.
\end{align*}
In comparing the numerical SH field with the function $A_1\cos(\theta-\phi)+A_3\cos(3(\theta-\phi))$, we find a maximum absolute error of $.023$ and a mean absolute error of $.009$. We provide plots of the recovered phase, wave number, amplitudes, phase gradients, sh field, approximated field, and field error.
\begin{figure}
\centering
\includegraphics[scale=0.4]{PhaseEst_v5_mu_03.png}
\includegraphics[scale=0.4]{WaveNums_v5_mu_03.png} \\[2ex]
\includegraphics[scale=0.4]{PhaseGradients_v5_mu_03.png}
\includegraphics[scale=0.4]{AmpEsts_v5_mu_03.png} \\[2ex]
\includegraphics[scale=0.3]{FieldEst_v5_mu_03.png}
\caption{knee bend fits}
\end{figure}

\subsection{Stationary Dislocation}
\par For the stationary dislocation, we choose to keep $k_{B}$ fixed to 1, and optimize for the parameters $\beta$, and $\phi$. We find more difficulty reconstructing a close fit to the SH field, though there are some refinements to the experiment that can be made. Specifically, we can integrate the SH field for a longer period of time, and modify the domain to include more stripes to sample regions further away from the defect. Nevertheless, we still see that the stationary dislocation phase solution provides a reasonable estimate of the SH field over most of the domain. Below, we display the objective function and gradient norm after 10000 iterations.
\begin{verbatim}
Step:  10000
Gradient Norm 0.0035862456230889967
Obj Function Norm:  0.004736415200037534
New Vals:  1.0 0.5531030652369063 -0.13494478559678233
Est Field max err: 0.5317768804964573
Est Field mean err: 0.03553398813306181
\end{verbatim}

\begin{figure}
\centering
\includegraphics[scale=0.7]{FieldEst_v1_0408_2.png}
\caption{Stationary Dislocation Fit}
\end{figure}
\FloatBarrier




\section{Experiments in Systems Identification}
\par A major aspiration of this research is to deduce order parameter equations from microscopic data. SH is an attractive model to use as a test case, as the wave vector dynamics are clearly a valid macroscopic description. While many succesful recoveries of microscopic models have been catalogued by the SINDy algorithm, successful identifications of models in terms of other coordinates have been far fewer. In this section, we discuss challenges related to identifying RCN like models from SH data.
\subsection{Plain SINDy for SH and PGB}
\par As a preliminary first step, we combined ideas from Kutz \cite{KUTZ2016} and Schaefer \cite{SCHA2017} to discover the Swift Hohenberg PDE from data consisting of solutions on the square.
\par Using a square grid of dimension $256 \times 256$ on the interval $[-25\pi,25\pi]^2$, we integrate Swift-Hohenberg from $t=0$ to $t=100$ using time step $h = .5$. Thus, we obtain a dataset of dimension $256 \times 256 \times 201$. We used backward finite differences to compute $200$ time derivatives $u_t(x_i,y_j,t_k)$. We then form various features out of the data $u(x_i,y_j,t_k)$. These features can be any term we suspect might be in the PDE, and primarily consists of polynomial functions, trigonometric functions, and spatial derivatives of $u$. We then form the matrix equation $u_t = \Theta(u)\xi$, where $\Theta$ is a matrix whose columns are the features constructed from $u$. The task then is to solve $u_t = \Theta(u)\xi$ where $\xi$ is a sparse vector of coefficients. Figure 4.1 shows an example of the matrix equation one might form where $u$ depends on only one spatial variable.
\begin{figure}
\centering
\includegraphics[scale=0.5]{pde_sindy_lib.png}
\caption{SINDy equation for PDEs with one spatial variable}
\end{figure}
\par Once the time derivatives are obtained and feature matrix is constructed, the task is to approximate a solution of
\begin{equation}
\xi = \text{argmin}_{\hat{\xi}}\| \Theta(u)\hat{\xi}-u_t\|_2^{2}+\lambda \|\hat{\xi}\|_{0}.
\end{equation}
In practice, what we actually solve is
\begin{equation}
\xi = \text{argmin}_{\hat{\xi}}\| \Theta(u)\hat{\xi}-u_t\|_2^2+\lambda \|\hat{\xi}\|_{1}.
\end{equation}
Following the approach of Kutz, we use Sequential Threshold Ridge Regression with a test/train split, outlined in algorithm 1 and algorithm 2.
\begin{figure}
\centering
\includegraphics[scale=0.5]{STRidge.png}
\caption{Algorithm 1: STRidge}
\end{figure}
\begin{figure}
\centering
\includegraphics[scale=0.5]{TrainSTRidge.png}
\caption{Algorithm 2: TrainSTRidge}
\end{figure}
\par Following the approach of Schaefer, we use coarsened grids to reduce the data being passed to the optimization procedure, and use spectral methods to form the spatial derivatives being passed to our library. Note that all polynomials, spatial derivatives and time derivatives are reshaped into column vectors and stacked vertically. For this preliminary proof of concept, our library is constructed with the terms $1,u,u^2,u^3,u^4,u_x, u_y, u_x^2, u_y^2, u_x u_y, u_{xy}, u_{xx}, u_{yy},\nabla^2u$, and $\nabla^4 u$. The recovered model is shown in figure 4.4, which agrees with Swift-Hohenberg.
\begin{figure}
\centering
\includegraphics[scale=0.8]{PDEFindCode2.png}
\caption{SH Model Recovery}
\end{figure}

\par We also used SINDy to validate the self-dual equation governing phase grain boundaries. Using the same phase we recovered in the PGB fit experiment, we computed numerical estimates of $\theta_{x}, \theta_{y}, \theta_{xx}, \theta_{yy}, |k|, \nabla \cdot \vec{k}, \nabla \times \vec{k}, \text{det}(J)$ from the recovered phase, and used functions of these terms to run a SINDy procedure on the phase surface, with the hope of recovering the self-dual equation
\begin{equation}
    (\nabla \cdot \vec{k})^2 -1 + 2|k|^2 - |k|^4 = 0.
\end{equation}
We test three left hand sides, $(\nabla \cdot \vec{k})^2, (\nabla \cdot \vec{k})^2 -1 + 2|k|^2 - |k|^4, $ and $2|k|^2 - |k|^4$, with corresponding candidate functions chosen arbitrarily.
We get the following fits:
\begin{verbatim}
u_t = (0.000636 +0.000000i)|k|
    + (-1.973117 +0.000000i)|k|^2
    + (1.000082 +0.000000i)|k|^4
    + (0.028843 +0.000000i)div(k)
    + (0.000233 +0.000000i)div(k)^4
    + (0.972378 +0.000000i)1
    + (0.000000 +0.000000i)theta_y
\end{verbatim}

\begin{verbatim}
u_t = (0.001469 +0.000000i)div(k)
\end{verbatim}

\begin{verbatim}
u_t = (0.001237 +0.000000i)div(k)
    + (-0.999817 +0.000000i)div(k)^2
    + (-0.002423 +0.000000i)|k|
    + (1.003394 +0.000000i)1
    + (2.734268 +0.000000i)Jk
    + (-0.000970 +0.000000i)e^(div(k))
    + (0.000000 +0.000000i)sin(theta_y)
\end{verbatim}

\FloatBarrier

\subsection{Simultaneous discover of coordinates and dynamics}
\par Ultimately, what one wishes to obtain is an algorithm that can deduce the approriate order parameter equations from the microscopic data, without any foreknowledge of the appropriate coordinates. As an example, would it be possible to extract phase and amplitude equations for the seed configurations of a sunflower head? The use of autoencoder neural networks have been demonstrated to recover dynamics in terms of an alternative coordinate system. Consider for example, a video of a swinging pendulum. Kutz et al were able to combine an autoencoder with the SINDy framework to obtain the equation $\ddot{\theta} = \sin(\theta)$ \cite{CHAMP2019}. However, we are not aware of examples in the literature where this method recovered order parameter equations from PDEs. Nevertheless, we are interested in regularizing/specializing this method for the problem at hand.
\par The method trains a neural network to generate its own input, with a latent space (typically of reduced dimension) in the middle. Thus, we can break the autoencoder into two components, the encoder and the decoder. If $x$ is our input data, $z$ the latent space representation of $x$, and $\hat{x}$ the output, the autoencoder can be represented as
\begin{equation}
\hat{x} = \psi(z) = \psi(\varphi(x)).
\end{equation}
In this case $\varphi$ is the encoder and $\psi$ is the decoder. The method consists of minimizing the autoencoder error, coupled with minimizing the SINDy error applied to the latent space representation. The loss function for a single example is then
\begin{equation}
L(x) = \|x - \psi(z)\|_2^2 + \lambda_1\|(\nabla_x z)\dot{x}-\Theta(z)\xi\|_2^2 + \lambda_2\|\xi\|_1.
\end{equation}
\par While there are some documented successes with this approach in the literature, making this work in the case of SH fields has proven to be quite challenging. Indeed, even the task of extracting the phase surface from SH data is an open problem. We have opted to break this task into incremental steps, which we document in this section. 
\subsubsection{Autoencoding SH Fields}
\par Due to the flexibility of neural networks, achieving an autoencoder of arbitrary latent dimension is easy, which we achieve using both a standard feed forward network and a convolutional network. We use data of dimension $1001 \times 128 \times 128$, and select our latent space to be of dimension $32 \times 32$. We run the autoencoders on 10 random samples from data that was not touched during training, and visualize the latent space, along with the input and its autoencoder approximation. Those comparisons are seen in figure 4.5 for the feedforward autoencoder, and 4.6 for the convolutional autoencoder.
\begin{figure}
\centering
\includegraphics[scale=0.35]{ae5.png}
\caption{Feedforward Autoencoder Latent Space}
\end{figure}
\begin{figure}
\centering
\includegraphics[scale=0.35]{ae6.png}
\caption{Convolutional Autoencoder Latent Space}
\end{figure}

\subsubsection{Non-ML to Estimate the Phase/Wave Vector}
\par While the long term goal is to have a framework that can machine learn the correct order parameter, we first considered methods of phase extraction based on first principles. These include a heuristic method, wavelet-like methods, and the use of the Hilbert Transform. 
\par The heuristic method is based on the intuitive notion of what the wave vector is. The key properties are that the wave vector should point in the direction of greatest change of the microscopic field, and its norm should capture the wavelength of the pattern in that direction. It is this second property that introduces some ambiguity, as the phase function will not be perfectly linear, but rather slowly varying. Formally, the precise meaning of the wave vector is as the phase gradient. However, away from defects, it seems reasonable to estimate the field with a linear phase, and use the resulting fit to define $\vec{k}$. The heuristic algorithm is based on this principle.
The microscopic field is denoted by $w$. The general procedure to obtain the wave vector has two stages.
\par In stage one, the wave vector direction is found at a point $r_0$ by creating a small disk of radius $\epsilon$ about $r_0$, $B_{\epsilon}(r_0)$. We select an optimal point on this disk $r_{\text{opt}}$ via
\[
    r_{\text{opt}} = \max_{r \in B_{\epsilon}(r_0)}|w(r)-w(r_0)|.
\]
By constrcution of the set $B_{\epsilon}(r_0)$, $r_{\text{opt}}$ is identified with an angle $\phi$, which we take to be the wave vector direction.
\par In stage two, we use the obtained directions to form profiles of the $w$ field. In the simplest scenario, the wave length at $(x_0,y_0)$ is found by moving in the direction $(\cos{\phi},\sin{\phi}))$ until a critical point is found. This distance is saved and denoted by $d_{\text{forward}}$. Then this is repeated, but in the direction $(-\cos{\phi},-\sin{\phi})$. This direction is saved and denoted by $d_{\text{backward}}$. The wavelength is taken to be $2(d_{\text{forward}}+d_{\text{backward}})$. If the point $(x_0,y_0)$ happens to be a min/max, then the procedure is the same, except we do not multiply the sum by $2$.
\par In practice, there are certain scenarios where the above procedure does not yield the correct wave length. Most notably, there are points in the knee bend where one should only traverse in one of the (forward or backward) directions. This occurs for points near the phase grain boundary. Of course, the RCN derivation is not intended to be valid near defects. However, even at the phase grain boundary, the phase is well defined, and its gradient exists. Thus, it is not unreasonable to recover the wave vector on the PGB. In fact, the wave vector is accurate for points on the PGB. It is the points just above and below the PGB where the algorithm breaks down. Figure 4.7 provides plots of the exact and recovered wave numbers, along with the absolute error.

\begin{figure}
\centering
\includegraphics[scale=0.3]{approx_wavenums_mu_03.png}
\includegraphics[scale=0.3]{approx_wavenums_mu_07.png}\\[2ex]
\includegraphics[scale=0.3]{exact_wavenums_mu_03.png}
\includegraphics[scale=0.3]{exact_wavenums_mu_07.png}\\[2ex]
\includegraphics[scale=0.3]{wavenum_err_mu_03.png}
\includegraphics[scale=0.3]{wavenum_err_mu_07.png}
\caption{Top Row: Approximate wavenumbers for shallow (left) and sharp (right) knee bends. Middle Row: Exact wavenumbers. Bottom Row: Absolute error}
\end{figure}

\par Another approach to estimating the wave vector is to use a Gaussian bump to yield a localized patch of the field, and find the dominant component of the spectrum. This is in the same spirit as a wavelet algorithm explored by Bowman \cite{BOWMAN1998250}. The results of this approach are depicted in Figure 4.8
\begin{figure}
\centering
\includegraphics[scale=0.4]{spec_wave_nums_mu_03.png}
\includegraphics[scale=0.4]{spec_wave_nums_mu_07.png}
\caption{Wavelet approach: Left column corresponds to shallow knee bend, right column is a sharp knee bend.}
\end{figure}
\FloatBarrier

\par A particularly interesting method investigated by Lega \cite{LEGA2018} makes use of the Hilbert Transform. Given a continuous, periodic signal, the Hilbert Transform returns a sensible estimate of its phase. In the case of the PGB and zipper defects, with defect core along the x-axis, we computed the hilbert transform of the horizontal profiles for every $y$. We then unwrap the profile, and use unwrapping again in the vertical direction to get a smooth phase. The resulting surface, while not precisely the phase, yields patterns that share the same tological character as the original pattern. This technique was used recently to demonstrate that the zipper configuration indeed has large phase gradients coming from the $y$ component of the wave vector. The results of this method are depicted in figure 4.9. 
\begin{figure}
\centering
\includegraphics[scale=0.7]{HilbertPhaseMu_03.png}\\[2ex]
\includegraphics[scale=0.9]{HilbertPhaseMu_07.png}
\caption{Hilbert Transform Approach}
\end{figure}

\par We note here that none of these methods provide a highly accurate phase surface. In the case of the heuristic approach, using a wavelength estimate based on distance to traverse one period doesn't account for nonlinearities in the phsae. In the case of the Hilbert Transform, one would ideally be able to consider signal profiles that are in the direction of the wave vector, which is very difficult numerically.
\FloatBarrier


\subsubsection{Phase Extraction via NN}
\par Having explored first principles methods to estimate the phase, we turn now to explore the use of a neural network to extract phases from fields of the form $w = \cos(\theta)$, where $\theta$ has no singularities, and is mostly slowly varying. To generate a rich dataset, we use as a base the phase surface for a phase grain boundary, and allow for rotations in the spatial coordinates. To encourage irregularities/distortions in the dataset, we use a Taylor Series approximation to $\text{cosh}(x)$, and use a gaussian sample to perturb its coefficients. The end result is a set of patterns mostly similar to what one may see in an SH field, though there are some frames where the wavenumber distribution is not as confined. Using $128 \times 128$ pattern frames, and there corresponding $32 \times 32$ coarsened phases, we train a simple feedforward neural network to extract the phase from the pattern. The result of this experiment is depicted in figure 4.10.

\begin{figure}[h]
\centering
\includegraphics[scale=0.35]{phase_extrat_nn_simple.png}
\caption{Phase Extract NN}
\end{figure}
\FloatBarrier


\subsubsection{Phase Extraction via NN Regularization}
\par While it is certainly promising to extract the phase using a standard feedforward network where the phase is used in the loss function, this ability is insufficient to resolve the problem at hand. We are after a simultaneous discovery of dynamics and order parameters, and therefore such a method will not have the benefit of a precomputed order paramater to penalize the loss. The next step in achieving such a task is to see if the order parameter can be learned by a neural net without supplying it as a penalty. Thus, one must consider using hard and soft constraints.
\par We are currently working on a simple experiment for such an approach, where we take as our dataset known solutions of the self-dual equation, ie knee bends of varying angles. We propose to use an autoencoder to reconstruct each frame, where the loss function includes a reconstruction loss, and a self-dual penalty applied to the latent variable. 
\begin{equation}
    L = \text{MSE}(w,\hat{w}) + \lambda \text{MSE}\left((\nabla \cdot \vec{k})^2-1+2|k| -|k|^2\right),
\end{equation}
where $\vec{k} = (\frac{\partial z}{\partial x},\frac{\partial z}{\partial y})$, $z$ being the latent variable. Let $w^{i}$ be the $i^{th}$ frame of an SH simulation. Then $z^{i} = \psi(w^{i})$ is the encoded state of $w^{i}$. One needs to compute $\Delta z^{i}$, which involves $2^{nd}$ order derivatives of the encoded variable with respect to the spatial coordinates. An interesting note, is that a priori there is no reason to believe that the encoded variable resides on a regular cartesian grid. It is the hope that this regularization allows us to make that interpretation. To compute such terms, we need to access the derivatives of the encoder with respect to the input $w^{i}$. ie,
\begin{equation*}
    \frac{\partial z}{\partial x} =  \frac{\partial \psi(w)}{\partial x} = \frac{\partial \psi}{\partial w}\frac{\partial w}{\partial x}.
\end{equation*}
Thus, the procedure must involve extracting derivatives of the encoder with respect to its input, as well as precomputing derivatives of the microscopic data with respect to its native coordinates. We attempted to perform such a regularization in tensorflow, but found the latent variable had no obvious relation to a phase surface. A careful code analysis must be conducted to check that the gradients of the encoder with respect to the expected input are being calculated correctly. 
\par Beyond the self-dual soft constraint, it is unclear how best to regularize the encoder to encourage spatial smoothness in the latent variable. A good first step could be to start the process by training a neural network that maps patches of SH-like fields to precomputed, coarsened phases, and use the weights of this network as the initial state of the encoder. More generally, the issue of spatial and temporal smoothness of the latent space is an open problem with many proposed regularizations.


\section{Next Steps}

\par The first open objective is to realize an RCN phase from the latent space of an autoencoder. The key challenges here are to understand how to utilize the computational graph in either pytorch/tensorflow, and determine if further regularization or pre-training is required to encourage the latent variable to be phase-like. The issue of smoothness of the latent variable has been explored recently \cite{XIE2024}, and the proposed methods should be explored and utilized.
\par Assuming a successful phase extraction in the form of a latent variable, we will be positioned to start working on the joint optimization procedure, where the order parameter and its dynamics are learned simultaneously. A review of the recently proposed standard \cite{FRIES2022} and weak \cite{TRAN2023} forms of simultaneous discovery should be useful here. As a side note, one of the main contributors to these methods is interested in collaborating with me on the challenge of identifying order parameter equations from SH fields.
\par Outside of the systems identification goal, there are still many interesting tasks that could be undertaken. In particular, variational analysis and numerical methods are lacking for the scenario of convection on the ellipse. If analysis could be done to determine the locations of the dislocations on the major axis of the ellipse, one would be able to prescribe a specific boundary condition and attempt to realize an RCN phase yielding the properties observed in the SH simulation. Analogously, one could attempt to setup a variational problem with appropriate boundary conditions to argue that RCN phases predict the defects in the ellipse.

\section{Code}
Code for all numerical simulations, machine learning experiments, and phase extraction algorithms are available in the following repositories: \newline
\href{https://github.com/EMcDugald/NaturalStripePatterns}{https://github.com/EMcDugald/NaturalStripePatterns} \newline
\href{https://github.com/EMcDugald/ConvectionPatterns}{https://github.com/EMcDugald/ConvectionPatterns}


\bibliographystyle{plain}
%\nocite{*}
\bibliography{refs}



\end{document}


